
\documentclass[11pt]{article}

\usepackage{amsmath}
\usepackage{amssymb}
\usepackage{amsthm}
\usepackage{indentfirst}
\usepackage{natbib}
\usepackage[colorlinks=true,allcolors=blue]{hyperref}
\usepackage{doi}
\usepackage{url}

\DeclareMathOperator{\iscontr}{isContr}
\DeclareMathOperator{\isprop}{isProp}
\DeclareMathOperator{\isset}{isSet}
\DeclareMathOperator{\idtype}{Id}

\newcommand{\universe}{\mathcal{U}}

\let\emptyset=\varnothing

\newtheorem{theorem}{Theorem}

\begin{document}

\title{More on Identity Types}

\author{Charles J. Geyer}

\maketitle

\section{Introduction}

This is about the ``argument'' (in scare quotes) that Roy Cook and I had
in the FIG (Foundations Interest Group of the Minnesota Center for Philosophy
of Science) discussion of \citet{geyer-survey}.

It seems I need to make a few points about identity types that I either omitted
from the discussed paper or did not emphasize enough.

\section{Mere Propositions}

A type $A$ is a \emph{proposition} if
\begin{equation} \label{eq:isProp}
  \isprop(A) :\equiv \prod_{x, y : A} \idtype_A(x, y) 
\end{equation}
\citep[Definition~3.3.1 in][hereinafter referred to as the
HoTT book]{hott-book}) \citep[Section~8.3]{geyer-survey}.
\defcitealias{hott-book}{HoTT book}
The \citetalias{hott-book} often uses \emph{mere proposition} to emphasize that
they are using this concept, but, strictly speaking, the ``mere'' is
superfluous.

We can read this as $A$ has at most one term if we are giving this the usual
reading in propositions-as-types, that is,
$$
   \forall_{x : A} \forall_{y : A} (x = y)
$$
But the homotopy-theoretic reading is that $A$ has at most one
\emph{path component:} it is either the empty space or a path connected space.

(Mere) propositions are not the lowest level of the homotopy level hierarchy.
Below is types satisfying
\begin{equation} \label{eq:isContr}
   \iscontr(A) :\equiv \sum_{x : A} \prod_{y : A} \idtype_A(x, y)
\end{equation}
the propositions-as-types reading for this being
$$
   \exists_{x : A} \forall_{y : A} (x = y)
$$
These are (the homotopy-theoretic reading) the \emph{contractible spaces}
or (the sets-as-types reading) the singleton sets.

Note an interesting curiosity, the ZFC hierarchy begins $\emptyset$,
$\{\emptyset\}$, $\ldots,$ where there is nothing special about the particular
singleton set chosen to be the second level of the hierarchy.
In contrast, the HoTT hierarchy \citepalias[Chapter~7]{hott-book}
\citep[Section~8.11]{geyer-survey} starts with contractible types (singleton
sets) and the next level is (mere) propositions (subsingleton sets).
So the hierarchies are sort of reversed.  But also the next level of the
HoTT hierarchy is (mere) sets, which includes the whole universe of
(HoTT) sets.

We emphasize the following, which was not emphasized enough in
\citet{geyer-survey}, which is Lemma~{3.3.2} in the \citetalias{hott-book}.
\begin{theorem} \label{th:prop}
If $A$ is an inhabited (true) (mere) proposition, then $A \simeq 1$.
\end{theorem}
That is, $A$ is (homotopy) equivalent to the type $1$.
Hence by univalence $A =_\universe 1$.

Of course, there are also (mere) propositions that are not inhabited (false).
Those are (homotopy) equivalent to the type $0$.

\section{How Identity Types Work}

\subsection{Contractible Types}

If $\iscontr(A)$, then $A \simeq 1$.  This follows from
Theorem~\ref{th:prop} because all contractible types are inhabited.

Hence $x, y : A$ implies $x =_A y$.  $A$ has exactly one point.

\subsection{(Mere) Propositions}

If $\isprop(A)$, then either $A \simeq 1$, in which case we say it is
true or $A \simeq 0$, in which case we say it is false.

Considered as sets, these are subsingleton sets (either singleton or empty).

\subsection{(Mere) Sets}

A type $A$ is a \emph{set} if
\begin{equation} \label{eq:isSet}
  \isset(A) :\equiv \prod_{x, y : A} \prod_{p, q : \idtype_A(x, y)}
  \idtype_{\idtype_A(x, y)}(p, q)
\end{equation}
\citepalias[Definition~{3.1.1} in the]{hott-book})
\citep[Section~8.3]{geyer-survey}.
I sometimes say \emph{mere set} to emphasize that
we are using this concept, but, strictly speaking, the ``mere'' is
superfluous.

So here again equality (identity types) behave as one intuits from classical
logic.  If $A$ is a (mere) set and $x, y : A$, then $x = y$ may be true
(inhabited) or false (not inhabited), but that is all that needs to be said.
Another way to say this is that $\idtype_A(x, y)$ is a mere proposition
so works the way one expects by analogy with classical logic.

Read homotopy-theoretically, of course, we do not think of the elements (terms)
of a (mere) set (type $A$ satisfying $\isset(A)$) as points in a space but
as path-connected components.  But that is just a change in language rather
than a change in the mathematics.

\subsection{Anima}

As \citet{geyer-survey} emphasizes a general type $A$ may be an \emph{anima}
(also called weak $\infty$-groupoid) in that it may have nontrivial paths,
paths between paths, paths between paths between paths, and so forth
ad infinitum.

This (something \citet{geyer-survey} did not emphasize) comes from
\emph{coinduction:} identity types are types.

But (another thing \citet{geyer-survey} did not emphasize and should have)
if we ever get to a level where an identity type is contractible, then there
is no higher-level nontrivial homotopy information.
\begin{theorem} \label{th:cut}
If $x, y : 1$, then $\idtype_1(x, y) \simeq 1$.
\end{theorem}
This is Theorem~{2.8.1} in the \citetalias{hott-book}.
It is an important contribution of Voevodsky.

The feeling that equalities are not behaving the way you intuit can stop
at any (homotopy) level.  Any time you get a contractible type, whether
this is the type you started with, or an identity type thereof, or an
identity type of an identity type, or an identity type of an identity
type of an identitiy type, or with arbitrarily many repetitions of 
an identity type of, then you are done.  There is nothing interesting
in any higher-level identity types.

\begin{thebibliography}{}

\bibitem[Geyer(2026)]{geyer-survey}
Geyer, C.~J. (2026),
\newblock On generalizations of equality, logic, and set theory in
    mathematics of the past century.
\newblock \url{https://github.com/cjgeyer/Survey}.
\newblock Read by the Foundations Interest Group of the Minnesota Center
    for Philosophy of Science October 1, 2026.

\bibitem[Hofmann and Streicher(1998)]{hofmann-streicher}
Hofmann, M., and Streicher, T. (1998).
\newblock The groupoid interpretation of type theory.
\newblock In G. Sambin and J.~M. Smith, (eds.),
    \emph{Twenty-five years of constructive type theory (Venice, 1995)}.
\newblock Oxford Logic Guides, \textbf{36}, 83--111.
\newblock Oxford University Press, New York.
\newblock \doi{10.1093/oso/9780198501275.003.0008}.

\bibitem[Martin-L\"{o}f(1975)]{martin-lof-mltt}
Martin-L\"{o}f, P. (1975).
\newblock An intuitionistic theory of types: predicative part.
\newblock In: H. E. Rose, J. C. Shepherdson (eds.),
    \emph{Logic Colloquium ‘73, Proceedings of the Logic Colloquium,
    Studies in Logic and the Foundations of Mathematics}, \textbf{80}, 73--118.
\newblock Elsevier.
\newblock \doi{10.1016/S0049-237X(08)71945-1}.

\bibitem[The Univalent Foundations Program(2013)]{hott-book}
The Univalent Foundations Program (2013).
\newblock \emph{Homotopy Type Theory: Univalent Foundations of Mathematics}.
\newblock Institute for Advanced Study, Princeton, NJ.
\newblock \url{https://homotopytypetheory.org/book}.

\end{thebibliography}

\end{document}

This about the mistake in the document about what one can prove constructively.

First we discuss triple negation reduction.  For all $P$ we have
$$
   \neg \neg \neg P \to \neg P
$$
is provable constructively
(for all $P$ not just $P$ that are mere propositions).
\citep[Exercise~1.11 in][hereinafter referred to as the HoTT book]{hott-book}.
\defcitealias{hott-book}{HoTT book}

This means if we think of the infinite sequence of statements
\begin{gather*}
   P
   \\
   \neg P
   \\
   \neg \neg P
   \\
   \neg \neg \neg P
   \\
   \neg \neg \neg \neg P
   \\
   \vdots
\end{gather*}
all but the first three logically imply one of the first three.

In fact, all but the first three are logically equivalent to the first three
because
$$
   P \to \neg \neg P
$$
is also provable constructively.

The question now arises whether for any $P$ we can prove something other
than
\begin{gather*}
   P
   \\
   \neg P
   \\
   \neg \neg P
\end{gather*}
The law of noncontradiction, also provable constructively
\citep[Section~1.1]{bauer-constructive}, says for all $P$
$$
   \neg ((\neg P) \times (\neg \neg P))
$$

And this is where I stumbled.  There is no proposition $P$ that is neither
true nor false, but that does not mean that we can constructively prove which
(that would be excluded middle if it were true for all $P$).
Similarly, there is no $P$ such that $\neg P$ is neither true nor false,
but that does not mean that we can constructively prove which
(that would be the erroneous statement asserted in Section 6
of \citet{geyer-survey} if it were true for all $P$).

